\chapter{The Diagonal}
% OUTLINE Ch2 "The Diagonal (the one gesture under everything)". Laws:
% both inherited (citation + intuition, labeled blocks adds-no-claim);
% non-funge continued. CITATION: Cantor 1891 (uncountability, the
% diagonal argument); Russell 1901/1902 (the class of classes); Richard
% 1905 + the Berry paradox (Russell/Berry ~1906) as the cautionary
% cousins told as the trade's own. The GESTURE named (apply a
% representation to itself, negate at the diagonal); the "paradoxes" =
% the gesture run in an unfenced language. Zero device appearances.
% Gate: Kodo.

\begin{intuition}
Someone claims to have listed every possible person's name-list --- every
list of names there could be, numbered, list one, list two, on forever.
Walk down the diagonal: take the first name off list one, the second off
list two, the $n$th off list $n$, and change each one. The name-list you
have just written differs from list one in its first entry, from list two
in its second, from every list somewhere --- so it is on none of them.
The claimant listed every list and you built one they missed, from
nothing but their own list turned against itself. The feeling to carry:
hand a describer its own description, twist the diagonal, and you have a
thing it cannot hold.
\end{intuition}

That maneuver is the single most productive gesture in this book, and it
has one author and one birth year for its clean form: Georg Cantor, 1891.
Cantor's target was sizes of infinity. He supposed the real numbers
between zero and one could be listed --- first, second, third, each
written as an endless decimal --- and then built a new decimal by walking
the diagonal and changing each digit: differ from the first number in the
first place, the second in the second, and so on. The result is a genuine
number between zero and one that appears nowhere on the supposed complete
list, so no such list is complete: the reals are a strictly larger
infinity than the counting numbers. What matters to this volume is not
the size result, celebrated as it is, but the \emph{shape} of the
argument, because that shape is about to appear in three more houses
wearing three more costumes. The shape: given any attempt to list all
objects of a kind by describing each, form the object that disagrees,
along the diagonal, with every description on the list --- and conclude
that the attempt failed.

Say the gesture in the volume's own terms, because naming it is the
chapter's work. A listing is a representation: a way of assigning
descriptions to things. The diagonal object is built by \emph{applying
the representation to itself} --- reading down the list's own entries ---
and then \emph{negating} at the point where the list describes its own
$n$th place. Self-application, then negation. That is the entire engine.
The quine of the last chapter performed the self-application and stopped
--- it printed its own text, no negation, no flip, a fixed point that
rests. Add the negation and the fixed point cannot rest: the object that
disagrees with itself along the diagonal has nowhere on the list to sit,
and the disagreement is the proof. The reader should hold the two halves
apart, because the whole book lives in their combination: self-reference
alone reproduces; self-reference plus negation refutes.

The gesture is so sharp that pointed carelessly it draws blood, and the
subject's early history is a cabinet of exactly that. Bertrand Russell,
in 1901, pointed the diagonal at sets themselves: consider the class of
all classes that do not contain themselves, and ask whether it contains
itself --- either answer forces its opposite. It is Cantor's diagonal
with sets for lists and membership for the digit, and it cut the ground
out from under the era's foundations. The same year's neighbors are a
rogues' gallery of the gesture unfenced: Jules Richard, in 1905, aimed it
at the phrases that define numbers and produced a defined number that
defies its own definition; the Berry sentence --- ``the least number not
nameable in fewer than twelve words,'' which has just named it in eleven
--- aims it at nameability. Each is the diagonal run in a language allowed
to talk about its own descriptions without restriction, and each returns
a contradiction. The trade learned to read these not as disasters but as
\emph{measurements}: a contradiction from the diagonal is the system
reporting that it let a representation range over itself unfenced. The
paradoxes are the gesture; the fence is what tells a theorem from a
catastrophe.

\begin{intuition}
The feeling for the two outcomes, side by side: the barber who shaves
exactly those who do not shave themselves cannot answer whether he shaves
himself --- run to contradiction, the unfenced version. But the same
self-application, fenced --- asking not for a paradoxical barber but for
what a consistent system can and cannot prove about itself --- returns
not a contradiction but a \emph{limit}, stated and survivable. Same
gesture; the fence decides whether it explodes or informs.
\end{intuition}

Two honesty notes, as the chapters before.

The non-funge law, sharpening now that negation is in play: the diagonal
object is built \emph{from a particular list}, and it is that list's
refuter, no other's. Cantor's missing real is missing from the specific
enumeration it diagonalizes; Russell's class is bound to the specific
comprehension it exploits. There is no universal diagonal object floating
free of a list --- the gesture always needs a particular representation
to turn against, and produces a residue bound to that representation. The
family resemblance across the cases is total; the objects are never the
same object. The book's coming chapters are, each one, this gesture in a
new house, and the whole point of the sixth chapter is that the houses
do not share their rooms.

And the mechanism note: the diagonal proves that certain listings are
incomplete or certain classes incoherent. It does not, of itself, say
anything about what can be known, thought, or physically done --- those
are the readings the preface fenced, and this chapter, which is pure
combinatorics of lists, supplies them no support whatever.

With self-application and negation now named as one engine, the book can
enter the house where the engine did its most famous work --- where the
list is the theorems of arithmetic, the diagonal sentence speaks of its
own provability, and the residue that survives is the one the trade calls
incompleteness. That is the next chapter, and the first of the four
great rooms.
